% Generated by Litex; source parsed, not verified.
% Compile with XeLaTeX.
\documentclass[11pt]{article}
\usepackage{amsmath,amssymb}
\usepackage{fontspec}
\setmainfont[BoldFont=cmunbx.otf,ItalicFont=cmunit.otf]{cmunrm.otf}
\usepackage[margin=25mm]{geometry}
\setlength{\parindent}{0pt}
\setlength{\parskip}{0.4em plus 0.1em minus 0.1em}
\setlength{\abovedisplayskip}{0.65em plus 0.15em minus 0.1em}
\setlength{\belowdisplayskip}{0.65em plus 0.15em minus 0.1em}
\setlength{\emergencystretch}{2em}
\raggedbottom
\begin{document}
\par\medskip\noindent\textbf{Claim.}\par\nopagebreak[4]
\par\noindent For every $n\in \mathbb{N}_{>0}$, we have $n = 1 \lor n > 1$.\par
\par\medskip\noindent\textbf{Proof.}\par\nopagebreak[4]
$n \geq 1$;
\par\noindent To prove $n = 1 \lor n > 1$, consider the following cases.\par
\begin{description}
\item[\textbf{Case 1.}] Suppose $n = 1$.\par
$n = 1 \lor n > 1$;
\item[\textbf{Case 2.}] Suppose $n \ne 1$.\par
\par\noindent We prove $n > 1$ by contradiction. Suppose $\neg\left(n > 1\right)$.\par
$n \leq 1$;
$n = 1$;
\par\noindent This contradicts $n \ne 1$.\par
\end{description}
\par\nobreak\hfill$\square$\par
\end{document}
